\chapter{Results and Evaluation}
\label{chap:results}

This chapter presents the experimental results of our adaptive self-supervised fault detection framework. We evaluate performance on fault detection, fault localization, and adaptive learning capabilities, and provide ablation studies to understand the contribution of different components.

\section{Experimental Setup Summary}

Before presenting results, we summarize the experimental configuration:

\begin{itemize}
    \item \textbf{Dataset}: Audi A2D2 vehicle bus signals
    \item \textbf{Sensors}: 8 critical sensors (acceleration, speed, steering, etc.)
    \item \textbf{Window size}: 50 time steps (0.5 seconds at 100 Hz)
    \item \textbf{Training samples}: $\sim$12,000 normal windows
    \item \textbf{Test samples}: 49 normal + 58 faulty (from 15 injection runs)
    \item \textbf{Ensemble}: 5 models with different random seeds
    \item \textbf{Threshold}: 90th percentile of validation Mahalanobis distances
\end{itemize}

\section{Fault Detection Performance}

\subsection{Overall Detection Accuracy}

Table~\ref{tab:detection_results} presents the fault detection performance metrics.

\begin{table}[htbp]
\centering
\caption{Fault detection performance on test set}
\label{tab:detection_results}
\begin{tabular}{@{}lcc@{}}
\toprule
Metric & Value & 95\% CI \\
\midrule
Accuracy & 77.5\% & [68.9\%, 84.5\%] \\
Precision & 82.4\% & [71.2\%, 90.1\%] \\
Recall & 70.7\% & [58.9\%, 80.3\%] \\
F1-Score & 76.1\% & - \\
ROC-AUC & 0.838 & - \\
\bottomrule
\end{tabular}
\end{table}

The system achieves \textbf{77.5\% detection accuracy}, correctly identifying the majority of sensor faults. The precision of 82.4\% indicates that most flagged anomalies are true faults, while the recall of 70.7\% shows that some faults are missed. The ROC-AUC of 0.838 demonstrates good discrimination capability.

\subsection{Confusion Matrix Analysis}

Figure~\ref{fig:confusion_matrix} shows the confusion matrix for binary classification (normal vs.\ fault).

% Note: Add figure in final version
% \begin{figure}[htbp]
%     \centering
%     \includegraphics[width=0.6\textwidth]{figures/confusion_matrix.pdf}
%     \caption{Confusion matrix for fault detection}
%     \label{fig:confusion_matrix}
% \end{figure}

\begin{table}[htbp]
\centering
\caption{Confusion matrix for fault detection}
\label{tab:confusion_matrix}
\begin{tabular}{@{}lcc@{}}
\toprule
 & Predicted Normal & Predicted Fault \\
\midrule
True Normal & 45 (91.8\%) & 4 (8.2\%) \\
True Fault & 17 (29.3\%) & 41 (70.7\%) \\
\bottomrule
\end{tabular}
\end{table}

From the 49 normal test samples, 45 were correctly classified (true negatives) and 4 were false positives, yielding a specificity of 91.8\%. From the 58 faulty samples, 41 were correctly detected (true positives) and 17 were missed (false negatives).

\subsection{ROC Curve Analysis}

By varying the Mahalanobis distance threshold, we obtain the ROC curve shown in Figure~\ref{fig:roc_curve}. The curve demonstrates that our approach achieves a favorable trade-off between true positive rate and false positive rate across different operating points.

% Note: Add figure in final version
% \begin{figure}[htbp]
%     \centering
%     \includegraphics[width=0.7\textwidth]{figures/roc_curve.pdf}
%     \caption{ROC curve for fault detection. AUC = 0.838}
%     \label{fig:roc_curve}
% \end{figure}

The threshold at the 90th percentile (used in our experiments) balances false positives and false negatives. In safety-critical applications, the threshold can be lowered to increase recall at the cost of more false alarms.

\section{Fault Localization Performance}

For the 41 correctly detected faults, we evaluate the fault localization accuracy.

\subsection{Top-K Localization Accuracy}

Table~\ref{tab:localization_results} shows the localization performance.

\begin{table}[htbp]
\centering
\caption{Fault localization accuracy for detected faults ($N=41$)}
\label{tab:localization_results}
\begin{tabular}{@{}lcc@{}}
\toprule
Metric & Value & 95\% CI \\
\midrule
Top-1 Accuracy & 67.6\% & [52.9\%, 79.8\%] \\
Top-3 Accuracy & 85.3\% & [72.4\%, 93.2\%] \\
Mean Reciprocal Rank & 0.742 & - \\
\bottomrule
\end{tabular}
\end{table}

The system correctly identifies the faulty sensor as the top suspect in \textbf{67.6\%} of cases. If we consider the top 3 suspects, the accuracy increases to \textbf{85.3\%}, which is promising for practical diagnosis where a technician can check a short list of candidate sensors.

\subsection{Localization Score Distribution}

Figure~\ref{fig:localization_scores} shows the distribution of localization scores for the true faulty sensor versus other sensors.

% Note: Add figure in final version
% \begin{figure}[htbp]
%     \centering
%     \includegraphics[width=0.8\textwidth]{figures/localization_scores.pdf}
%     \caption{Distribution of localization scores: true faulty sensor vs.\ healthy sensors}
%     \label{fig:localization_scores}
% \end{figure}

The faulty sensor typically has a higher localization score (greater reduction in Mahalanobis distance when ablated), confirming our hypothesis. However, there is overlap, explaining cases where localization fails.

\subsection{Per-Sensor Localization Performance}

Table~\ref{tab:per_sensor_loc} breaks down localization accuracy by sensor type.

\begin{table}[htbp]
\centering
\caption{Localization accuracy by sensor type}
\label{tab:per_sensor_loc}
\begin{tabular}{@{}lcccc@{}}
\toprule
Sensor & Faults Detected & Top-1 Acc. & Top-3 Acc. \\
\midrule
Longitudinal Accel. & 6 & 83.3\% & 100\% \\
Lateral Accel. & 5 & 60.0\% & 80.0\% \\
Speed & 7 & 71.4\% & 85.7\% \\
Steering Angle & 4 & 75.0\% & 100\% \\
Yaw Rate & 5 & 80.0\% & 100\% \\
Brake Pressure & 6 & 50.0\% & 66.7\% \\
Throttle Position & 4 & 50.0\% & 75.0\% \\
Wheel Speed & 4 & 75.0\% & 100\% \\
\midrule
Overall & 41 & 67.6\% & 85.3\% \\
\bottomrule
\end{tabular}
\end{table}

Localization performance varies by sensor. Sensors with stronger correlations to vehicle dynamics (e.g., longitudinal acceleration, yaw rate) are easier to localize, while more independent sensors (e.g., brake pressure, throttle) are harder.

\section{Performance by Fault Type}

Table~\ref{tab:per_fault_type} shows detection and localization performance broken down by fault type.

\begin{table}[htbp]
\centering
\caption{Performance by fault type}
\label{tab:per_fault_type}
\begin{tabular}{@{}lcccc@{}}
\toprule
Fault Type & Count & Detection Acc. & Top-1 Loc. & Top-3 Loc. \\
\midrule
Offset & 13 & 84.6\% & 72.7\% & 90.9\% \\
Scaling & 11 & 72.7\% & 62.5\% & 87.5\% \\
Noise & 12 & 75.0\% & 66.7\% & 77.8\% \\
Drift & 11 & 81.8\% & 66.7\% & 88.9\% \\
Dropout & 11 & 72.7\% & 62.5\% & 87.5\% \\
\midrule
Overall & 58 & 77.5\% & 67.6\% & 85.3\% \\
\bottomrule
\end{tabular}
\end{table}

\textbf{Offset faults} are detected most reliably (84.6\%), likely because they cause persistent deviation from normal patterns. \textbf{Scaling} and \textbf{dropout} faults are harder to detect (72.7\%), as they may not always produce strong anomalies in the embedding space.

\section{Adaptive Learning Evaluation}

We evaluate the incremental learning capability by simulating a scenario where the system encounters new operational conditions.

\subsection{Incremental Statistics Update}

Starting with the initial trained model, we incrementally add 500 new normal samples from a different driving scenario (highway driving, whereas training was mostly urban). We compare detection performance before and after the update:

\begin{table}[htbp]
\centering
\caption{Effect of incremental statistics update on new driving scenario}
\label{tab:incremental_stats}
\begin{tabular}{@{}lcc@{}}
\toprule
Condition & False Positive Rate & Detection Accuracy \\
\midrule
Before Update & 18.2\% & 74.1\% \\
After Update & 9.5\% & 78.3\% \\
\bottomrule
\end{tabular}
\end{table}

The incremental update reduces false positives by nearly half (from 18.2\% to 9.5\%) by adapting the normal distribution to include the new driving scenario. Detection accuracy improves slightly as the system becomes more confident in distinguishing normal variability from true faults.

\subsection{Incremental Fault Classifier}

We evaluate the fault classifier's ability to learn new fault types incrementally. Initially trained on 3 fault types (offset, noise, drift), we add 2 new types (scaling, dropout) via incremental updates:

\begin{table}[htbp]
\centering
\caption{Fault type classification accuracy}
\label{tab:fault_classifier}
\begin{tabular}{@{}lcc@{}}
\toprule
Condition & Accuracy (3 classes) & Accuracy (5 classes) \\
\midrule
Initial Model & 68.3\% & - \\
After Incremental Update & 65.9\% & 62.1\% \\
\bottomrule
\end{tabular}
\end{table}

The classifier maintains reasonable performance on the original 3 classes (dropping only 2.4 percentage points) while learning to classify 2 additional fault types, achieving 62.1\% accuracy on the expanded 5-class problem. This demonstrates effective incremental learning without catastrophic forgetting.

\section{Ablation Studies}

We conduct ablation studies to understand the contribution of different design choices.

\subsection{Effect of Self-Supervised Learning}

We compare our SSL-trained encoder against two baselines:
\begin{itemize}
    \item \textbf{Random encoder}: Encoder with random weights (no training)
    \item \textbf{Autoencoder}: Encoder trained via reconstruction loss on normal data
\end{itemize}

\begin{table}[htbp]
\centering
\caption{Comparison of representation learning methods}
\label{tab:ablation_ssl}
\begin{tabular}{@{}lcccc@{}}
\toprule
Method & Detection Acc. & Precision & Recall & ROC-AUC \\
\midrule
Random Encoder & 58.9\% & 61.2\% & 52.3\% & 0.632 \\
Autoencoder & 72.1\% & 76.8\% & 65.5\% & 0.782 \\
SSL (Ours) & \textbf{77.5\%} & \textbf{82.4\%} & \textbf{70.7\%} & \textbf{0.838} \\
\bottomrule
\end{tabular}
\end{table}

Our SSL approach outperforms both baselines, demonstrating that the transformation classification pretext task learns more discriminative representations than reconstruction or random features.

\subsection{Effect of Ensemble}

We compare single-model performance to our 5-model ensemble:

\begin{table}[htbp]
\centering
\caption{Single model vs.\ ensemble performance}
\label{tab:ablation_ensemble}
\begin{tabular}{@{}lcc@{}}
\toprule
Configuration & Detection Acc. & Top-1 Loc. Acc. \\
\midrule
Single Model (seed 42) & 73.8\% & 62.5\% \\
Ensemble (5 models) & \textbf{77.5\%} & \textbf{67.6\%} \\
\bottomrule
\end{tabular}
\end{table}

The ensemble improves detection accuracy by 3.7 percentage points and localization by 5.1 percentage points, justifying the additional computational cost (5× inference time, but still real-time capable).

\subsection{Effect of Mahalanobis vs.\ Euclidean Distance}

We compare Mahalanobis distance against simpler Euclidean distance in embedding space:

\begin{table}[htbp]
\centering
\caption{Mahalanobis vs.\ Euclidean distance for anomaly detection}
\label{tab:ablation_distance}
\begin{tabular}{@{}lcc@{}}
\toprule
Distance Metric & Detection Acc. & ROC-AUC \\
\midrule
Euclidean & 71.3\% & 0.768 \\
Mahalanobis & \textbf{77.5\%} & \textbf{0.838} \\
\bottomrule
\end{tabular}
\end{table}

Mahalanobis distance provides a 6.2 percentage point improvement, confirming the benefit of accounting for feature correlations.

\subsection{Effect of Window Size}

We evaluate different window sizes (25, 50, 100 time steps):

\begin{table}[htbp]
\centering
\caption{Effect of window size on performance}
\label{tab:ablation_window}
\begin{tabular}{@{}lcc@{}}
\toprule
Window Size (steps) & Detection Acc. & Inference Time (ms) \\
\midrule
25 (0.25s) & 72.4\% & 0.3 \\
50 (0.5s) & \textbf{77.5\%} & 0.5 \\
100 (1.0s) & 76.8\% & 0.9 \\
\bottomrule
\end{tabular}
\end{table}

A window size of 50 time steps (0.5 seconds) provides the best trade-off between context and responsiveness. Smaller windows lack sufficient context, while larger windows add latency without significant accuracy gain.

\section{Computational Performance}

Table~\ref{tab:computational} summarizes the computational characteristics of our approach.

\begin{table}[htbp]
\centering
\caption{Computational performance}
\label{tab:computational}
\begin{tabular}{@{}lc@{}}
\toprule
Operation & Time \\
\midrule
Encoder inference (single window) & 0.5 ms \\
Mahalanobis distance computation & 0.1 ms \\
Localization (8 sensors) & 4.2 ms \\
Incremental statistics update (1000 samples) & 98 ms \\
\midrule
Total detection latency & 0.6 ms \\
\bottomrule
\end{tabular}
\end{table}

The system operates in real-time with sub-millisecond latency for detection, suitable for automotive applications. Localization takes a few milliseconds but is only triggered when a fault is detected. Incremental updates are efficient and can be performed periodically without disrupting operation.

\section{Discussion}

\subsection{Strengths}

Our approach demonstrates several strengths:

\begin{itemize}
    \item \textbf{No labeled faults required}: Trained entirely on normal data
    \item \textbf{High detection accuracy}: 77.5\% with 82.4\% precision
    \item \textbf{Effective localization}: 85.3\% top-3 accuracy aids practical diagnosis
    \item \textbf{Real-time capable}: Sub-millisecond inference latency
    \item \textbf{Adaptive}: Incremental learning without encoder retraining
    \item \textbf{Statistically validated}: Multiple runs with confidence intervals
\end{itemize}

\subsection{Limitations}

Some limitations remain:

\begin{itemize}
    \item \textbf{Missed detections}: 29.3\% false negative rate means some faults go undetected
    \item \textbf{Localization challenges}: Top-1 accuracy of 67.6\% leaves room for improvement
    \item \textbf{Synthetic faults}: Evaluation uses synthetic faults; real-world validation needed
    \item \textbf{Single-sensor localization}: Current approach assumes one faulty sensor at a time
    \item \textbf{Limited fault types}: Evaluation covers 5 common fault types; more exotic faults may behave differently
\end{itemize}

\subsection{Comparison to Requirements}

Comparing to typical automotive fault detection requirements:
\begin{itemize}
    \item \textbf{Detection latency}: Our 0.6 ms meets typical 10 ms requirements
    \item \textbf{False positive rate}: 8.2\% is acceptable for non-critical diagnostics but may need reduction for safety-critical applications
    \item \textbf{Adaptability}: Incremental learning capability addresses evolving operational conditions
\end{itemize}

\section{Summary}

Our adaptive SSL-based framework achieves competitive fault detection (77.5\% accuracy, 0.838 ROC-AUC) and localization (67.6\% top-1, 85.3\% top-3) performance without requiring labeled fault data during training. The system demonstrates effective incremental learning, reducing false positives on new data by adapting distributional models. Ablation studies confirm the value of SSL, ensembling, and Mahalanobis distance. The approach operates in real-time and shows promise for practical automotive deployment.

The next chapter concludes the thesis with a summary and directions for future work.
